% !PRINTED 434
\weberpage{434}
%<<CONT
consequently the second term of (20) is divisible by $p^2$, and if
one substitutes $\gamma p^2$ for it, it follows that
\[
\alpha\delta - \beta\gamma = 1.
\]
We thus arrive at the result:

IV. \gesperrt{The invariant equation $F_p(v,u)$ has multiple roots always
then and only then, when $u=j(\omega)$, where $\omega$
is the root of a quadratic equation whose discriminant
is not divisible by $p$, but which admits a solution of equation (18).}

For a given $p$ there is only a finite number of
values $m$ which meet this requirement, since $m < 4p^2$ must
hold. The values $m=3$, $m=4$ are contained among them as special
cases.

Since by (18) $-m$ is a quadratic residue of $p$, $p$
is representable by a form of discriminant $-m$; but $p^2$ is
representable by the principal form of a discriminant. The
class of forms by which $p$ is representable is therefore itself either
the principal class, or, compounded with itself, it gives the
principal class; it is therefore two-sided. We can therefore also
state Theorem IV. in the following way:

V. \gesperrt{The roots of the discriminant of the invariant equation
$F_p(v,u)$ are the singular invariants $j(\omega)$, where $\omega$
is the root of a quadratic equation of negative discriminant
$D$ for which the prime $p$ has the index 1 or 2.}

\begin{center}\rule{0.18\linewidth}{0.4pt}\end{center}
